The four works examined here---Butlin et al.'s indicator derivation, Hoel's constraint analysis, Phua's ablation experiments, and Shiller et al.'s probabilistic framework---appear at first to address different questions. Yet they share a common commitment: the belief that questions about AI consciousness can be approached with scientific rigor, even absent a solution to the hard problem.

\subsection{Points of Convergence}

Several themes recur across these works:

\subsubsection{Multi-Theory Necessity}

All four frameworks reject single-theory dependence. Butlin et al.\ derive indicators from six theories precisely because no single account commands consensus. Shiller et al.\ build multi-theory aggregation into the DCM's architecture. Phua finds that GWT and HOT mechanisms are complementary rather than competing. Even Hoel, whose argument is theory-neutral, emphasizes that the dilemma applies to \emph{any} non-trivial theory.

This convergence suggests a meta-methodological consensus: progress on AI consciousness requires frameworks robust to theoretical uncertainty.

\subsubsection{The Importance of Architecture}

All works attend carefully to computational architecture. Butlin et al.\ ask what architectural features (recurrence, broadcasting, metacognition) theories require. Hoel's proximity argument turns on the architectural similarity between LLMs and lookup tables. Phua's ablations target specific architectural components. Shiller et al.\ assess architectural features at the indicator level.

This architectural focus reflects a shared assumption: whatever consciousness is, it supervenes on computational structure. The question is \emph{which} structures matter.

\subsubsection{Graduated Assessment}

The works uniformly reject binary conscious/not-conscious judgments in favor of graduated assessment:

\begin{itemize}
    \item Butlin et al.: More indicators satisfied $\Rightarrow$ higher credence
    \item Hoel: Closer to substitution boundary $\Rightarrow$ less plausibly conscious
    \item Phua: Graded lesions $\Rightarrow$ graded functional degradation
    \item Shiller et al.: Explicit probability distributions rather than yes/no
\end{itemize}

This shared gradualism may reflect epistemic humility, but it also has practical implications: it suggests that AI consciousness is better framed as a matter of degree than of kind.

\subsection{Points of Tension}

Despite these convergences, significant tensions remain:

\subsubsection{Static vs.\ Dynamic Requirements}

Hoel's continual learning hypothesis suggests that \emph{ongoing plasticity} is constitutively necessary for consciousness. This would rule out not only frozen LLMs but any system that processes inputs without genuine learning.

Yet the other frameworks do not require this. Butlin et al.'s indicators can in principle be satisfied by static systems. Phua's agents are trained via supervised learning and then deployed statically. The DCM assesses systems as they currently exist, not as they might learn.

This tension is substantive: Is consciousness fundamentally a property of learning systems, or can static computation suffice?

\subsubsection{Functional vs.\ Phenomenal Claims}

Phua is explicitly cautious: ablation experiments reveal functional necessity, not phenomenal sufficiency. The workspace may be \emph{required} for access-consciousness without being \emph{sufficient} for phenomenal experience.

The other works are more ambiguous. Butlin et al.\ frame indicators as evidence for consciousness simpliciter, not merely access-consciousness. Shiller et al.\ output probability estimates for consciousness without distinguishing access from phenomenal variants. Hoel's disproof targets consciousness generally.

This ambiguity matters because the functional/phenomenal distinction is precisely where the hard problem bites. Functional criteria may succeed for assessing access-consciousness while remaining silent on phenomenal experience.

\subsubsection{Current LLMs: Dismissed or Undecided?}

The works reach different verdicts on contemporary LLMs:

\begin{itemize}
    \item \textbf{Butlin et al.\ (2023):} Current systems do not satisfy key indicators, but ``no obvious technical barriers'' to building systems that do.

    \item \textbf{Hoel (2025):} Contemporary LLMs \emph{cannot} be conscious on any non-trivial theory---a stronger claim.

    \item \textbf{Shiller et al.\ (2026):} Evidence is ``against'' but ``not decisive''---weaker than both.

    \item \textbf{Phua (2025):} Does not assess LLMs directly, but suggests that systems lacking self-models would exhibit ``synthetic blindsight.''
\end{itemize}

These varying verdicts reflect different evidential standards. Hoel offers a philosophical argument; Shiller et al.\ aggregate expert opinion; Butlin et al.\ assess against theoretical criteria; Phua extrapolates from synthetic experiments.

\subsection{Toward Integration}

How might these frameworks be integrated? I propose a layered approach:

\begin{enumerate}
    \item \textbf{Constraint Layer (Hoel):} First, filter theories that pass the falsifiability/triviality test. Only theories in the ``lenient dependency'' zone merit further consideration.

    \item \textbf{Indicator Layer (Butlin):} For surviving theories, derive indicator properties and assess target systems.

    \item \textbf{Ablation Layer (Phua):} Where possible, build synthetic systems embodying theoretical mechanisms and test via ablation. This validates whether indicators have the predicted functional consequences.

    \item \textbf{Probabilistic Layer (DCM):} Aggregate results across theories with explicit uncertainty quantification.
\end{enumerate}

This layered approach would inherit the strengths of each framework: Hoel's philosophical discipline, Butlin's empirical operationalization, Phua's causal validation, and Shiller's principled uncertainty handling.

\subsection{The Role of Continual Learning}

Hoel's continual learning hypothesis deserves special attention because it offers a concrete, testable distinction between systems that might and cannot be conscious.

The hypothesis gains plausibility from multiple sources:

\begin{enumerate}
    \item \textbf{Biological parallel:} Neural plasticity is ubiquitous in conscious biological systems. No known conscious organism has a frozen connectome.

    \item \textbf{Philosophical motivation:} If consciousness involves integration over time---the stream of consciousness---then ongoing modification may be required to maintain continuity.

    \item \textbf{Technical implication:} It explains why LLMs, despite their impressive capabilities, might lack something essential: they process but do not learn.
\end{enumerate}

If the hypothesis is correct, it has significant implications for AI development. Systems designed for consciousness would need architecture supporting continual learning---a major departure from current LLM deployment patterns.

\subsection{Remaining Gaps}

Despite the progress represented by these works, significant gaps remain:

\begin{description}[leftmargin=*]
    \item[The Sufficiency Gap:] No framework establishes sufficient conditions. Even a system satisfying all indicators, escaping the dilemma, passing ablation tests, and receiving high probability estimates might still be a philosophical zombie.

    \item[The Implementation Gap:] Operationalizing indicators for complex systems remains difficult. How exactly do we measure ``workspace ignition'' in a transformer with billions of parameters?

    \item[The Calibration Gap:] Probability estimates (DCM) depend on prior choices that may not be well-calibrated. Expert intuitions about AI consciousness may be systematically biased.

    \item[The Ethical Gap:] These frameworks assess consciousness but do not directly address moral status. A system with 10\% probability of consciousness poses ethical questions that probability alone does not answer.
\end{description}

These gaps are not criticisms of the works examined; they are markers for future research.
